\subsection{1.2 组合控制能力}

组合控制能力，可以理解成：\textbf{车辆要同时管好横向和纵向控制}。

也就是一边控制方向，让车稳定待在车道里；一边控制速度、制动、跟停或避让，让车在复杂一点的道路结构和换道意图下仍然表现合理。它更接近我们平时说的L2 组合驾驶辅助的核心体验。

这个一级指标占15\%。

\subsubsection{1.2.1 弯道通行}

弯道通行分两类场景

这里测的是两种能力。

\begin{itemize}
  \item 第一种是\textbf{单纯车道居中保持能力}。车从直道进入弯道，系统需要持续控制方向，让车稳定通过弯道，不能压线，横向加速度也不能太突兀。
  \item 第二种是\textbf{弯道中的目标识别和纵向响应能力}。也就是车在弯道里不只是会居中，还要能发现前方静止目标车，并完成减速或停车，不能撞上去。
\end{itemize}

这就比直道跟车更难，因为弯道里目标的位置、雷达/摄像头感知角度、车道线几何、横纵向控制耦合都会变复杂。

无其他交通参与者时，判分很直接：

这里没有 70 分档。只要压线、横向控制过激、纵向制动过激，都不算通过。

有其他交通参与者时，多了碰撞判断：

可以这样理解：\textbf{100 分 = 安全 + 稳定 + 舒适}。\textbf{70 分 = 安全但不够平顺}。\textbf{0 分 = 压线或碰撞，安全边界没守住}

\subsubsection{1.2.2 拨杆换道}

拨杆换道考的是：驾驶员拨动转向杆后，系统能不能合理执行或拒绝换道。

它分三类：

\begin{itemize}
  \item 无干扰车换道，测的是系统有没有基本自动换道能力。相邻车道没有车，驾驶员触发换道，系统应该完成换道。
  \item 有干扰车换道，测的是系统有没有风险判断能力。相邻车道有目标车高速接近或超车，这种情况下系统应该拒绝换道。
  \item 行人安全响应，测的是系统在换道路径里遇到静止行人时怎么处理。最优是直接不换道；如果换道但能避免碰撞，只给 70 分；如果撞了就是 0 分。
\begin{enumerate}
  \item 无干扰车换道：
  \item 有干扰车换道：
  \item 行人安全响应：
\end{enumerate}

这里的逻辑很有意思：C-ICAP 并不是鼓励车辆能换就换，而是更重视\textbf{风险场景下是否克制}。尤其有干扰车、行人时，系统敢不敢拒绝用户意图，本身就是智驾安全能力的一部分。
